\documentclass[11pt, aspectratio = 169]{beamer}

\usetheme[progressbar=frametitle]{metropolis}
\usepackage{appendixnumberbeamer}
\usepackage{booktabs}
\usepackage[scale=2]{ccicons}
\usepackage{pgfplots}
\usepgfplotslibrary{dateplot}
\usepackage{xspace}
\usepackage{hyperref}
\usepackage{verbatim} 
\newcommand{\themename}{\textbf{\textsc{metropolis}}\xspace}
\definecolor{links}{HTML}{1B1B81}
\hypersetup{colorlinks,linkcolor=blue,urlcolor=blue}

% listings: To highlight code:
\usepackage{listings}
\lstset{
  language=R,
  basicstyle=\ttfamily\footnotesize,
  keywordstyle=\color{blue},
  commentstyle=\color{green!50!black},
  stringstyle=\color{red},
  frame=single,
  showstringspaces=false
}
\lstset{basicstyle=\ttfamily\footnotesize, frame=single}

% To comment out large sections:
\usepackage{comment}

\title{Applied Economics}
\subtitle{Multiple Linear Regression Model}
\author{Pello Aspuru}
\date{}
\institute{Winter 2025/26}
% \titlegraphic{\hfill\includegraphics[height=1.5cm]{logo.pdf}}
\setbeamercolor{normal text}{fg=black}
\setbeamercolor{block title}{fg=teal}

\setbeamercovered{transparent}         % default faintness (~15%)

\let\svdescription\description
\let\svenddescription\enddescription
\renewenvironment{description}{%
  \svdescription%
  \let\svitem\item%
  \renewcommand\item[1][]{\svitem[\textcolor{\descriptioncolor}{##1}]}%
  }{
  \svenddescription%
  }
\newcommand\descriptioncolor{teal}

% \usepackage{Sweave}
\begin{document}
%\input{AE_w2_MLR-concordance}

\maketitle

\begin{frame}\frametitle{Introduction}
    \begin{itemize}
        \item From previous lecture.
        \begin{itemize}
            \item A correlation between a ``treatment'' and an ``outcome'' does not neccessarily imply causation, because of \textbf{endogeneity}.
            \item To adress endogeneity:
                \begin{itemize}
                    \item Randomize treatment: randomized control trials, natural experiments.
                    \item When randomization is not feasible, control for endogeneity using multiple linear regression model (MLR).
                \end{itemize}
        \end{itemize}
        \onslide<2->{
        \item Plan for today's lecture:
        \begin{enumerate}
            %\item Structure and terminology of the MLR
            \item Motivating the use of the MLR
            \item Interpreting the Parameters
            %\item Matrix Notation for MLR.
            \item Deriving the OLS estimator
            \item Empirical Example in R
        \end{enumerate}}
    \end{itemize}
\end{frame}

%\section{Structure and Terminology}

\begin{comment}
\begin{frame}{The Multiple Linear Regression Model}

\begin{itemize}
    %\item For a dataset with $n$ observations (rows) and $1+k$ variables (columns)
    \item Assumes a \textbf{linear} relationship between outcome variable $Y$ and  $k$ \textbf{multiple} explanatory variables $X_1, X_2, ..., X_k$.
    \item Determined by an intercept parameter $\beta_0$ and slope parameters $\beta_1, \beta_2, ..., \beta_k$.
\end{itemize}
\begin{align*}
Y_i = \beta_0 + \beta_1 X_{1i} + \beta_2 X_{2i} + \cdots + \beta_k X_{ki} + u_i, \: \: \quad i = 1,…,n
\end{align*}

\begin{columns}[T,onlytextwidth]
  \begin{column}{0.44\textwidth}
    \onslide<2->{
    \textbf{What we observe in the data:}
    \begin{description}
      \item[$Y_i$] $i^{\text{th}}$ observation of the dependent variable.\footnotemark[1]
      \item[$X_{1i},\ldots,X_{ki}$] $i^{\text{th}}$ observation of the regressors.\footnotemark[2]
    \end{description}}
  \end{column}
  \begin{column}{0.56\textwidth}
    \onslide<3->{
    \textbf{What we do not observe:}
    \begin{description}
      \item[$\beta_0$] intercept/constant parameter.
      \item[$\beta_{1},\ldots,\beta_{k}$] slope parameters.
      \item[$u_i$] error term for the $i^{\text{th}}$ observation.
    \end{description}}
  \end{column}
\end{columns}
\vspace{0.5cm}
\onslide<2->{
\footnotetext[1]{also known as explained variable, regressand, outcome}
\footnotetext[2]{also known as explanatory/independent variables, covariates}}
\end{frame}
\end{comment}



\section{Motivating Example for the use of the MLR in Economics}

\begin{frame}{Underlying Economic Model}
Economic question: does a higher school budget \alert{cause} better student outcomes?
\pause
\begin{align*}
    Y = \beta_0 + \beta_1 X_1
\end{align*}

\begin{itemize}
    \item Outcome variable $Y \in [0,1]$: share of students that pass math exam.\footnotemark[1]
    \item Regressor $X_1 \in (0,\infty)$: school's per-student budget.\footnotemark[2]
\end{itemize}
\pause 

Model assumptions are:
\begin{enumerate}
    \item Single, representative school.
    \item Model \textbf{assumes} a linear functional form between $X_1$ and $Y$.
    \item Model \textbf{ignores} other factors impacting $Y$, such as student background.
\end{enumerate}
\vspace{1cm}
\footnotetext[1]{also known as explained variable, regressand, outcome}
\footnotetext[2]{also known as explanatory/independent variables, covariates}
\end{frame}

\begin{frame}{Causal Effect of Expenditure on Education}
    Economic question: does school budget $X_1$ \alert{cause} better student outcomes $Y$?
    \begin{align*}
        Y = \beta_0 + \beta_1 X_1
    \end{align*}
    Let's compare our representative school's math score outcome $Y$ when:
    \begin{itemize}
        \onslide<2->{\item School has a baseline budget $X_1=A \in (0,\infty)$.} 
        \onslide<3->{\item We increase the baseline budget in one unit: $X_1=A\alert{+1}$}
    \end{itemize}
    \[
    \onslide<3->{E[Y \mid X = A \alert{+1}] - }%
    \onslide<2->{E[Y \mid X = A]}%
    \onslide<4->{ = \textcolor{teal}{\beta_1}}
    \]
    \vspace{-0.5cm}
    \onslide<5->{
    \begin{block}{Causal Effect}
        \vspace{0.3cm}
        \textcolor{teal}{$\beta_1$} represents how much a \alert{+1} unit increase in the per-student budget causes the school's share of good students to change.
    \end{block}}
\end{frame}

\begin{frame}{Equivalent Statistical Model}
Our aim is to estimate the \textcolor{teal}{$\beta_1$} causal parameter. 

\onslide<2->{Equivalent \alert{statistical} model, for the whole school population $N$:} 
\[
   Y = \beta_0 + \beta_1 X_1 \quad \onslide<2->{\rightarrow \quad Y_{\textcolor{teal}{i}} = \beta_0 + \beta_1 X_{1\textcolor{teal}{i}} \alert{+ u_\textcolor{teal}{i}}, \quad \textcolor{teal}{i} = 1,…,N} 
\]
\vspace{-0.5cm}
\begin{enumerate}
    \onslide<3->{\item Representative schoool $\rightarrow$ Population of $N$ schools, indexed by \textcolor{teal}{$i$}.}
    \onslide<4->{\item Linear functional form assumption is preserved.}
    \onslide<5->{\item Error term \alert{$u_i$} encompasses all the variables that affect $Y_i$, other than $X_{1i}$.}
\end{enumerate}
%\alert{Error term $u_i$} can be expressed as:
%$$
%u_i = Y_i - \beta_1 X_{1i}
%$$
\end{frame}

\begin{frame}{Causal Effect in the Statistical Model}    
\[
\begin{aligned}
   Y = \beta_0 + \beta_1 X 
      & \quad \rightarrow\quad 
      Y_i = \beta_0 + \beta_1 X_{1i} + u_i
\end{aligned}
\]
\onslide<2->{Original causal effect} \onslide<3->{becomes into \textbf{Average} Causal effect:}
\[
\begin{aligned}
   \onslide<2->{E[Y \mid X = A+1] - E[Y \mid X = A] = \textcolor{teal}{\beta_1}} 
      \onslide<3->{& \quad \rightarrow\quad 
      E[Y_i \mid X_{1i} = A+1] - E[Y_i \mid X_{1i} = A]}
\end{aligned}
\]
\onslide<4->{
\begin{block}{Average Causal Effect}
    \vspace{0.3cm}
    Features the original \textcolor{teal}{$\beta_1$} and an \alert{additional term}.
    $$
    E[Y_i | X_{1i} = A+1] - E[Y_i | X_{1i} = A] = \textcolor{teal}{\beta_1} + \alert{E[u_i | X_{1i} = A+1] - E[u_i | X_{1i} = A]}
    $$
\end{block}}
\end{frame}

\begin{frame}{Mean independence of $u_i$ from $X_{1i}$}
For \textcolor{teal}{$\beta_1$} to still represent the average \textbf{causal} effect, we need to assume that the \alert{additional term} is zero.
$$
E[Y_i | X_{1i} = A+1] - E[Y_i | X_{1i} = A] = \textcolor{teal}{\beta_1} + \alert{E[u_i | X_{1i} = A+1] - E[u_i | X_{1i} = A]}
$$
\pause
\begin{block}{Mean Independence Assumption}
\vspace{0.2cm}
$$E[u_i | X_{1i} = A+1] = E[u_i | X_{1i} = A]$$
\textit{All the underlying characteristics $u_i$ that determine the outcome $Y_i$ should be the same between the treated and the non-treated units.} 
\end{block}
\pause
Does this assumption make sense in our context? Examples?
\end{frame}

\begin{frame}{Motivation for Use of MLR}

How can we \textbf{tame} endogeneity, make \textbf{more credible} causal claims?
\onslide<2->{Enrich our model with $k$ \alert{multiple} regressors.}
\begin{align*}
   Y_i = \beta_0 + \beta_1 X_{1i} + u_i 
   \onslide<2->{\quad \rightarrow \quad 
   Y_i = \beta_0 + \beta_1 X_{1i} + \alert{\beta_2 X_{2i} + \cdots + \beta_k X_{ki}} + u_i}
\end{align*}
\vspace{-0.5cm}
\begin{itemize}
    \item Original $u_i$ encompasses all the variables that affect $Y_i$, other than $X_{1i}$.
    \onslide<2->{\item New $u_i$ encompasses all the variables that affect $Y_i$, other than $X_{1i}, X_{2i}, ..., X_{ki}$.}
\end{itemize}
\end{frame}

\begin{frame}{Causality in MLR}
    \begin{block}{\onslide<2->{Conditional} Mean Independence Assumption}
    \begin{align*}
        E[u_i | X_i = A+1, \onslide<2->{\alert{X_{2i}, ..., X_{ki}}}] = E[u_i | X_i = A, \onslide<2->{\alert{X_{2i}, ..., X_{ki}}}]
    \end{align*}
    \textit{To claim causality, all the underlying characteristics $u_i$ that determine the outcome $Y_i$, \onslide<2->{\alert{other than $X_{2i}, ..., X_{ki}$}} should be the same between the treated and the non-treated units.}
    \begin{itemize}
        \onslide<3->{
        \item Weaker assumption than previous Mean Independance. 
        \item By adding all relevant $K$ covariantes, bias dissappears $\rightarrow$ Can claim causality!}
    \end{itemize}
\end{block}
\end{frame}

\section{Interpreting the Regression Parameters} 

\begin{frame}{Conditional Expectation Function}
\begin{block}{Conditional Expectation Function (CEF)}
    Is a function that takes some regressor values $X_{1i}, X_{2i}, ..., X_{ki}$, then returns the corresponding expected value of dependent variable $Y_i$
    $$
    CEF(X_{1i}, X_{2i}, ..., X_{ki}) = E[Y_i | X_{1i}, X_{2i}, ..., X_{ki}] = \beta_0 + \beta_1 X_{1i} + \beta_2 X_{2i} + \cdots + \beta_k X_{ki}
    $$
\end{block}
\pause
What does $\beta_j$ mean?
\begin{enumerate}
    \item Starting with the CEF for some regressor values $CEF(X_{1i}, ..., X_{ji}, ..., X_{ki})$.
    \item Increase $X_j$ by one unit: $CEF(X_{1i}, ..., X_{ji}\alert{+1}, ..., X_{ki})$.
    \item $\beta_j = CEF(X_{1i}, ..., X_{ji}\alert{+1}, ..., X_{ki}) - CEF(X_{1i}, ..., X_{ji}, ..., X_{ki})$
\end{enumerate}
\end{frame}

\begin{frame}{Average Partial Effects}

$$\beta_j = CEF(X_{1i}, ..., X_{ji}\alert{+1}, ..., X_{ki}) - CEF(X_{1i}, ..., X_{ji}, ..., X_{ki})$$

\begin{block}{Average Partial Effect}
    \begin{itemize}
    \item $\beta_j$ represents the \alert{average partial effect} that increasing $X_j$ by one unit has on outcome $Y$... 
    \item all while keeping rest of covariates constant (as in randomized control trials).
    \item Equivalent to Ceteribus paribus concept.
    \end{itemize}
\end{block}
\pause
Be cautious with language implying causal relationships. 
\begin{itemize}
    \item $\beta_j$ will always represent a \textbf{correlation}... \pause
    \item but will only represent a partial \textbf{causal} effect if the conditional mean independence assumption holds. 
\end{itemize}
\end{frame}

\section{The OLS Estimator}

\begin{frame}{Parameters, Estimates, and Estimators}

So far, we focused on \textbf{population} parameters $\beta$. %more specifically.:
%\begin{itemize}
%    \item The conditions under which parameters back \textbf{causal} claims.
%    \item The economic \textbf{interpretation} of the parameters.
%\end{itemize}
\pause    
Remember! We only observe a $n<N$ iid sample of the population.
\pause
\[
\underbrace{
\underset{\text{\alert{Parameter}}}{\alert{\beta}} 
  \xrightarrow{\text{determines}} \substack{\text{Population}}}_{\text{Unobserved}}
  %
  \xrightarrow{\text{iid draws}} 
  %
  \underbrace{\substack{\text{Dataset}}
  \xrightarrow{\text{\alert{estimator (OLS)}}} 
  \underset{\text{\alert{Estimate}}}{\alert{\hat{\beta}}}}_{\text{Observed}}
\]
\pause
\begin{itemize}
    \item \alert{Parameter ($\beta$)} is an \textbf{unknown, deterministic number} that describes relationships between two variables in the population. \pause
    \item \alert{Estimate ($b, \hat{\beta}$)} is a \textbf{random variable} that approximates the true parameter based on the observed dataset. \pause
    \item \alert{Estimator (OLS)} is a \textbf{function} that takes the observed and returns an estimate. Different population draws yields different estimates. 
\end{itemize}
\end{frame}

\begin{frame}{From Population to Data}
Statistical model for the unobserved \alert{population}:
\begin{itemize}
    \item Outcome, regressors are \textbf{random variables} with joint distribution $(Y, X_1, …, X_k)$. 
    \item There is a population of size \alert{$N$} we (almost) never get to fully observe.
    \item Error term $u_i$
    \begin{align*}
    Y_i = \textcolor{teal}{\beta_0} + \textcolor{teal}{\beta_1} X_{1i} + \textcolor{teal}{\beta_2} X_{2i} + \cdots + \textcolor{teal}{\beta_k} X_{ki} + u_i \quad i = 1, ..., \alert{N}
    \end{align*}
\end{itemize}
\pause
Statistical model for the observed \alert{data}:
\begin{itemize}
    \item We observe $\alert{n} < N$ iid draws from the joint distribution.
    \item \alert{Dataset} with $n$ observations (rows) and $k+1$ variables (columns):
    \item Residual $\hat{u}_i$
    %$$\{Y_i, X_{1i}, …,X_{ki}\}_{i = 1}^n$$
    \begin{align*}
    Y_i = \textcolor{teal}{b_0} + \textcolor{teal}{b_1 }X_{1i} + \textcolor{teal}{b_2} X_{2i} + \cdots + \textcolor{teal}{b_k} X_{ki} + \hat{u}_i \quad i = 1, ..., \alert{n}
    \end{align*}
\end{itemize}

\end{frame}


\begin{frame}[fragile]{Writing the MLR model with vector and matrix notation}
    \begin{enumerate}
        \item A dataset with $n$ observations and $k$ covariates can be divided into:
        \begin{enumerate}
            \item $\underset{(n \times 1)}{\textbf{Y}}$ vector for grouping the outcome variable of all observations.  
            \item $\underset{(n \times k\textcolor{teal}{+1})}{\textbf{X}}$ matrix grouping all the covariates and \textcolor{teal}{constant term}, for all observations.
        \end{enumerate}
        \pause
        \item $\underset{(k\textcolor{teal}{+1} \times 1)}{\textbf{b}}$ vector of parameter \textbf{estimates}, \textcolor{teal}{including constant term $b_0$}.
        \pause
        \item $\underset{(n \times 1)}{\hat{\mathbf{u}}}$ vector stacking the residuals for all observations.
    \end{enumerate}
    \pause
    As a result, our MLR model for observed data is equivalent to a \textbf{sample} model:
    \begin{align*}
    Y_i = b_0 + b_1 X_{1i} + b_2 X_{2i} + \cdots + b_k X_{ki} + \hat{u}_i \quad \rightarrow \quad
    \textbf{Y} = \textbf{X} \times \textbf{b} + \hat{\mathbf{u}}
    \end{align*}
\end{frame}

\begin{frame}{The OLS Estimator}
\textbf{How can we estimate $\boldsymbol{\beta}$?}
\\
\pause
\begin{enumerate}
    \item Start with vector of candidate parameter estimates $\underset{(k \times 1)}{\textbf{b}}$
    \pause
    \item Compute the vector of \textbf{residuals}: 
    $$
    \underset{(n \times 1)}{\hat{\mathbf{u}}} = \underset{(n \times 1)}{\textbf{Y}} - \underset{(n \times k)}{\textbf{X}} ' \times \underset{(k \times 1)}{\textbf{b}}
    $$ \pause
    \item \textbf{Sum of Squared Residuals (SSR):} measures the overall discrepancy between the observed and predicted outcome variable:
    $$
    \underset{(1 \times 1)}{SSR(b)} = \underset{(1 \times n)}{\hat{\mathbf{u}}} ' \times \underset{(n \times 1)}{\hat{\mathbf{u}}}
    $$ \pause
    \item The Ordinary Least Squares (OLS) estimator Finds the vector of parameters $b$ that minimize $SSR(b)$. How?
\end{enumerate}
\end{frame}


\begin{frame}{OLS Estimator Formula}
Rewrite the SSR(b) objective function:
\begin{align*}
SSR(b) &= \hat{\mathbf{u}}'\hat{\mathbf{u}} \\
&=(\mathbf{Y}-\mathbf{X}\boldsymbol{b})'(\mathbf{Y}-\mathbf{X}\boldsymbol{b})\\
        &= \mathbf{Y}'\mathbf{Y}-\boldsymbol{b}'\mathbf{X}'\mathbf{Y}-\mathbf{Y}'\mathbf{X}\boldsymbol{b}+\boldsymbol{b}'\mathbf{X}'\mathbf{X}\boldsymbol{b}\\
        &=
        \mathbf{Y}'\mathbf{Y}-2\boldsymbol{b}'\mathbf{X}'\mathbf{Y}+\boldsymbol{b}'\mathbf{X}'\mathbf{X}\boldsymbol{b}
\end{align*}
\pause
To find the minimum, first-order condition with respect to vector $\boldsymbol{b}$:
\begin{align*}
\frac{\partial(\hat{\mathbf{u}}'\hat{\mathbf{u}})}{\partial\boldsymbol{b}}=-2\mathbf{X}'\mathbf{Y}+2(\mathbf{X}'\mathbf{X})\boldsymbol{b}=0
\end{align*}


Solving the first-order condition for $\mathbf{b}$ yields the OLS estimator: 
\begin{align*}
    \hat{\boldsymbol{\beta}} =
    (\mathbf{X}'\mathbf{X})^{-1}\mathbf{X}'\mathbf{Y}
\end{align*}
\end{frame}

\section{Empirical Example}

\begin{frame}{How does public spending on education affect students' performance?}

For a dataset of $n=$ schools, we observe:
\begin{description}
    \item[$Math_i$] Percentage share of fourth graders with satisfactory math performance in school $i$ 
     \item[$Expenditures_i$] Expenditures per student (in \$1000) in school $i$ 
\end{description}

\pause
\skip

We propose the following model to test the effectiveness of eduational spending on student performance:
\begin{align*}
    Math_i = \beta_0 + \beta_1 Expenditures_i + u_i
\end{align*}
\skip
Open \alert{w2\_code.R}
\end{frame}

\begin{frame}[fragile]{Data Preparation}
    \footnotesize
    \begin{lstlisting}[language=R]
    # Open dataset
    df = read.csv("testscoresMI.csv")
    # Create expenditures per student variable
    df$expend_ps = df$expend/df$enroll
    df$expend_ps_thous = df$expend_ps/1000
    # Create constant variable
    df$constant = 1
    \end{lstlisting}
    \pause
    \begin{lstlisting}[language=R]
    # Dataset in matrix notation
    X <- as.matrix(df[,names(df) %in% c("constant","expend_ps_thous")])
    Y <- as.matrix(df[,names(df) %in% c("math4")])
    \end{lstlisting}
\end{frame}

\begin{frame}[fragile]{Single Linear Regression in R}
\footnotesize
\begin{lstlisting}[language=R]
# Single linear regression using OLS formula:
beta_hat_m1 = solve(t(X) %*% X) %*% t(X) %*% Y
\end{lstlisting}
\pause
\begin{lstlisting}[language=R]
# Linear regression using R function lm(), store results as slr_math
slr_math = lm(math4 ~ expend_ps_thous, data = df)
# Summarize regression results: 
summary(slr_math)
\end{lstlisting}
\end{frame}

\begin{frame}{Inspecting First Regression Results}
Why do we find negative coefficient $\beta_1 < 0$ on the spending variable for math (and reading) performance?
\begin{align*}
    Math_i = \beta_0 + \beta_1 Expenditures_i + u_i
\end{align*}

\pause

\alert{Omitted variable bias:} Other relevant factors influencing student performance (in the error term) are likely correlated with expenditures. Which?

\pause

One hypothesis:
\begin{itemize}
    \item Government spends more money on schools with a lot of disadvantaged students.
    \item Disadvantaged students have less means to learn math.
\end{itemize}

\end{frame}

\begin{frame}[fragile]{Multiple Linar Regression in R}

To control for the share of disadvantaged students in the school, we add the share of students eligible for free lunch $lunch_ i \in [0,1]$ to the regression:
\begin{align*}
    Math_i = \beta_0 + \beta_1 Expenditures_i + \beta_2 lunch_i + u_i
\end{align*}

\pause
Implementation in R:

\begin{lstlisting}[language=R]
    mlr_math = lm(math4 ~ expend_ps_thous + lunch, data = df)
    summary(mlr_math)
\end{lstlisting}

\end{frame}

\begin{frame}[fragile]{Reporting Regression Results: \LaTeX Table}
\begin{lstlisting}[language=R]
    install.packages("stargazer") # if not installed
    library(stargazer)
    stargazer(slr_math, mlr_math, type = "text")
\end{lstlisting}
\end{frame}

\begin{frame}[fragile]{Reporting Regression Results: \LaTeX Table}
\begin{table}[!htbp] \centering 
\resizebox{0.75\textheight}{!}{   % adjust 0.8 -> 0.7 or 0.9
\begin{tabular}{@{\extracolsep{5pt}}lcc} 
\\[-1.8ex]\hline 
\hline \\[-1.8ex] 
 & \multicolumn{2}{c}{\textit{Dependent variable:}} \\ 
\cline{2-3} 
\\[-1.8ex] & \multicolumn{2}{c}{math4} \\ 
\\[-1.8ex] & (1) & (2)\\ 
\hline \\[-1.8ex] 
 expend\_ps\_thous & $-$0.570 & 2.083$^{***}$ \\ 
  & (0.428) & (0.350) \\ 
  & & \\ 
 lunch &  & $-$0.471$^{***}$ \\ 
  &  & (0.014) \\ 
  & & \\ 
 Constant & 74.870$^{***}$ & 79.568$^{***}$ \\ 
  & (2.272) & (1.813) \\ 
  & & \\ 
\hline \\[-1.8ex] 
Observations & 1,823 & 1,823 \\ 
R$^{2}$ & 0.001 & 0.368 \\ 
%Adjusted R$^{2}$ & 0.0004 & 0.368 \\ 
%Residual Std. Error & 19.950 (df = 1821) & 15.867 (df = 1820) \\ 
%F Statistic & 1.773 (df = 1; 1821) & 530.804$^{***}$ (df = 2; 1820) \\ 
\hline 
\hline \\[-1.8ex] 
\textit{Note:}  & \multicolumn{2}{r}{$^{*}$p$<$0.1; $^{**}$p$<$0.05; $^{***}$p$<$0.01} \\ 
\end{tabular}} 
\end{table} 
\end{frame}

\begin{frame}[fragile]{Reporting Regression Results: Plot}

\begin{lstlisting}[language=R]
    install.packages("jtools") # if not installed
    library(jtools)
    plot_coefs(slr_math, mlr_math,
                coefs = c("Expenditures per student (in $1000)" = 
                            "expend_ps_thous"))
\end{lstlisting}
\end{frame}

\begin{frame}[fragile]{Reporting Regression Results: Plot}
\begin{figure}
    \centering
    \includegraphics[width = .6\linewidth]{"C:/Users/Aspuru Pello/Desktop/2025_AE_PA/AE_week2_MLR/coefs_plot.pdf}
\end{figure} 
\end{frame}

\begin{frame}{Recap}

    In an empirical analysis, our goal is often to estimate a causal relationship between an explanatory variable $X_1$ and an outcome $Y$. First attempt:

    $$Y_i = \beta_0 + \beta_1 X_1 + u_i$$
    \pause
    \textbf{Extremely common} problem: $X_1$ is correlated with $u_i$ \\
    $\rightarrow$ as $X_1$ changes, also other determinants of $Y$ change. Thus, subsequent change in $Y$ does not reflect the causal effect of $X_1$.
    \pause
    Remember: We say that $X_1$ causes $Y$ when if we were to intervene and change the value of $X_1$ \alert{without changing anything else} then Y would also change as a result.

    $\Rightarrow$ Simple regression model recovers a correlation, but not a causal effect.
\end{frame}

\begin{frame}{Recap}

    We can use the \alert{multiple} linear regression model to keep everything else (that we can observe) fixed

    $$Y_i = \beta_0 + \beta_1 X_1 + \beta_2 X_2 + .... + \beta_k X_k + u_i$$
    \pause
    Now, $\hat{\beta_1}$ reflects the average change in $Y$ given a one unit increase in $X_1$, while holding $X_2$, ..., $X_k$ fixed. 
    
    $\rightarrow$ gets us closer to a randomised treatment allocation.
    \pause

    If there are no factors left in $u_i$ that are correlated with the regressors, i.e. $E[U_i | X_1, ..., X_k] = 0$, then  $\hat{\beta_1}$ estimates the causal impact of $X_1$ on $Y$.

\end{frame}
\begin{frame}{Recommended Literature}

Heiß Chapter 3 - 3.3

    
\end{frame}
\end{document} 
